\documentclass[11pt,a4paper]{article}
% preamble.tex — estilo común de todos los apuntes Froth.
% Cada apunte hace % preamble.tex — estilo común de todos los apuntes Froth.
% Cada apunte hace % preamble.tex — estilo común de todos los apuntes Froth.
% Cada apunte hace \input{preamble} justo después de \documentclass.
\usepackage[margin=2.4cm]{geometry}
\usepackage{xcolor}
\definecolor{litblue}{RGB}{31,79,121}
\definecolor{litgreen}{RGB}{27,94,32}
\definecolor{litorange}{RGB}{230,81,0}
\usepackage{parskip}
\usepackage{titlesec}
\titleformat{\section}{\Large\bfseries\color{litblue}}{\thesection}{0.6em}{}
\titleformat{\subsection}{\large\bfseries\color{litblue}}{\thesubsection}{0.6em}{}
\usepackage{tcolorbox}
\usepackage{fancyhdr}
\pagestyle{fancy}
\fancyhf{}
\fancyhead[L]{\small\color{litblue}Froth — Apuntes de ML}
\fancyhead[R]{\small\thepage}
\renewcommand{\headrulewidth}{0.4pt}
\usepackage[colorlinks=true,linkcolor=litblue,urlcolor=litblue]{hyperref}

% Cabecera de cada apunte: \cabecera{numero}{titulo}
\newcommand{\cabecera}[2]{%
  \begin{tcolorbox}[colback=litblue,coltext=white,colframe=litblue,arc=2mm]
    {\Large\bfseries Apunte #1 — #2}\\[3pt]
    {\small Proyecto Froth · aprender Machine Learning construyendo}
  \end{tcolorbox}\vspace{0.8em}}

% Cajas temáticas
\newtcolorbox{concepto}[1]{colback=blue!4,colframe=litblue,title={Concepto: #1},arc=1.5mm}
\newtcolorbox{practica}{colback=green!5,colframe=litgreen,title={Pruébalo tú (teclado en mano)},arc=1.5mm}
\newtcolorbox{ojo}{colback=orange!8,colframe=litorange,title={Ojo / error típico},arc=1.5mm}
\newtcolorbox{resumen}{colback=gray!8,colframe=gray!60,title={En una frase},arc=1.5mm}

\newcommand{\codigo}[1]{\texttt{#1}}
 justo después de \documentclass.
\usepackage[margin=2.4cm]{geometry}
\usepackage{xcolor}
\definecolor{litblue}{RGB}{31,79,121}
\definecolor{litgreen}{RGB}{27,94,32}
\definecolor{litorange}{RGB}{230,81,0}
\usepackage{parskip}
\usepackage{titlesec}
\titleformat{\section}{\Large\bfseries\color{litblue}}{\thesection}{0.6em}{}
\titleformat{\subsection}{\large\bfseries\color{litblue}}{\thesubsection}{0.6em}{}
\usepackage{tcolorbox}
\usepackage{fancyhdr}
\pagestyle{fancy}
\fancyhf{}
\fancyhead[L]{\small\color{litblue}Froth — Apuntes de ML}
\fancyhead[R]{\small\thepage}
\renewcommand{\headrulewidth}{0.4pt}
\usepackage[colorlinks=true,linkcolor=litblue,urlcolor=litblue]{hyperref}

% Cabecera de cada apunte: \cabecera{numero}{titulo}
\newcommand{\cabecera}[2]{%
  \begin{tcolorbox}[colback=litblue,coltext=white,colframe=litblue,arc=2mm]
    {\Large\bfseries Apunte #1 — #2}\\[3pt]
    {\small Proyecto Froth · aprender Machine Learning construyendo}
  \end{tcolorbox}\vspace{0.8em}}

% Cajas temáticas
\newtcolorbox{concepto}[1]{colback=blue!4,colframe=litblue,title={Concepto: #1},arc=1.5mm}
\newtcolorbox{practica}{colback=green!5,colframe=litgreen,title={Pruébalo tú (teclado en mano)},arc=1.5mm}
\newtcolorbox{ojo}{colback=orange!8,colframe=litorange,title={Ojo / error típico},arc=1.5mm}
\newtcolorbox{resumen}{colback=gray!8,colframe=gray!60,title={En una frase},arc=1.5mm}

\newcommand{\codigo}[1]{\texttt{#1}}
 justo después de \documentclass.
\usepackage[margin=2.4cm]{geometry}
\usepackage{xcolor}
\definecolor{litblue}{RGB}{31,79,121}
\definecolor{litgreen}{RGB}{27,94,32}
\definecolor{litorange}{RGB}{230,81,0}
\usepackage{parskip}
\usepackage{titlesec}
\titleformat{\section}{\Large\bfseries\color{litblue}}{\thesection}{0.6em}{}
\titleformat{\subsection}{\large\bfseries\color{litblue}}{\thesubsection}{0.6em}{}
\usepackage{tcolorbox}
\usepackage{fancyhdr}
\pagestyle{fancy}
\fancyhf{}
\fancyhead[L]{\small\color{litblue}Froth — Apuntes de ML}
\fancyhead[R]{\small\thepage}
\renewcommand{\headrulewidth}{0.4pt}
\usepackage[colorlinks=true,linkcolor=litblue,urlcolor=litblue]{hyperref}

% Cabecera de cada apunte: \cabecera{numero}{titulo}
\newcommand{\cabecera}[2]{%
  \begin{tcolorbox}[colback=litblue,coltext=white,colframe=litblue,arc=2mm]
    {\Large\bfseries Apunte #1 — #2}\\[3pt]
    {\small Proyecto Froth · aprender Machine Learning construyendo}
  \end{tcolorbox}\vspace{0.8em}}

% Cajas temáticas
\newtcolorbox{concepto}[1]{colback=blue!4,colframe=litblue,title={Concepto: #1},arc=1.5mm}
\newtcolorbox{practica}{colback=green!5,colframe=litgreen,title={Pruébalo tú (teclado en mano)},arc=1.5mm}
\newtcolorbox{ojo}{colback=orange!8,colframe=litorange,title={Ojo / error típico},arc=1.5mm}
\newtcolorbox{resumen}{colback=gray!8,colframe=gray!60,title={En una frase},arc=1.5mm}

\newcommand{\codigo}[1]{\texttt{#1}}

\begin{document}
\cabecera{01}{Python, el motor; y el entorno virtual, la caja de herramientas}

\section{Python: qué es de verdad}

Los archivos \codigo{.py} de la carpeta \codigo{froth/} son texto plano: recetas escritas.
\textbf{Python} es el programa que sabe leerlas y ejecutarlas (el ``motor''). Por eso el
primer paso del proyecto fue instalarlo: sin motor, las recetas son papel mojado.

En tu máquina quedó instalado \textbf{Python 3.14}. Cuando escribes
\codigo{python archivo.py} en una terminal, le estás diciendo: ``motor, ejecuta esta receta''.

\begin{concepto}{librería (package)}
Casi nada se programa desde cero. Una \textbf{librería} es código empaquetado que otros
escribieron y tú instalas y reutilizas: \codigo{requests} (hablar con internet),
\codigo{pandas} (tablas de datos), etc. Se instalan con la herramienta \codigo{pip},
que viene con Python. Piensa en apps que le instalas al motor.
\end{concepto}

\section{El problema que resuelve el entorno virtual}

Imagina dos proyectos en tu PC: uno necesita \codigo{pandas} versión 2 y otro la versión 3.
Si instalas todo ``en el sistema'', se pisan entre sí y algo se rompe. La solución estándar:

\begin{concepto}{entorno virtual (venv)}
Una \textbf{carpeta-burbuja} (en nuestro caso \codigo{.venv/} dentro del proyecto) que
contiene una copia privada de Python y sus librerías, exclusiva de este proyecto.
Lo que instalas dentro no toca nada fuera, y viceversa.

Se creó con: \codigo{python -m venv .venv}

Y desde entonces, para ejecutar cualquier cosa del proyecto usamos el Python \emph{de la
burbuja}: \codigo{.venv\textbackslash Scripts\textbackslash python.exe}. Esa ruta rara del
principio de cada comando significa exactamente eso: ``usa el motor de ESTE proyecto''.
\end{concepto}

\section{requirements.txt: la lista de la compra}

El archivo \codigo{requirements.txt} lista las librerías que el proyecto necesita, con un
comentario de para qué sirve cada una. Cualquiera (tú en otra PC, un compañero) puede
recrear el entorno con una sola orden:
\begin{center}
\codigo{.venv\textbackslash Scripts\textbackslash python.exe -m pip install -r requirements.txt}
\end{center}
En Froth las instalamos \emph{por fases}: hoy solo las de la Fase 1 (\codigo{requests},
\codigo{pandas}, \codigo{pyarrow}, \codigo{truststore}). Las pesadas de ML llegarán cuando
toquen.

\begin{ojo}
Si un comando te dice \codigo{ModuleNotFoundError: No module named 'pandas'}, casi siempre
significa que ejecutaste con el Python \textbf{del sistema} en vez del de la burbuja.
Revisa que el comando empiece por \codigo{.venv\textbackslash Scripts\textbackslash python.exe}.
\end{ojo}

\begin{practica}
Abre PowerShell (tecla Windows $\rightarrow$ escribe \codigo{powershell} $\rightarrow$ Enter) y:
\begin{enumerate}
  \item \codigo{cd "C:\textbackslash Users\textbackslash you\textbackslash Documents\textbackslash Froth"}
  \item \codigo{.venv\textbackslash Scripts\textbackslash python.exe --version} \quad (debe decir 3.14)
  \item \codigo{.venv\textbackslash Scripts\textbackslash python.exe -m pip list} \quad (las librerías de tu burbuja)
\end{enumerate}
\end{practica}

\begin{resumen}
Python ejecuta el código; el \codigo{.venv} es la burbuja privada del proyecto con sus
librerías; \codigo{requirements.txt} es la lista para reconstruirla donde sea.
\end{resumen}

\end{document}
